\einheitenkopf{Einheit 4 von 5 · Gleichung bestimmen}
\zweigzeile{Hier lernst du, die Gleichung einer Geraden aus dem Graphen, aus Steigung und Punkt oder aus zwei Punkten zu bestimmen und den Schnittpunkt zweier Geraden zu berechnen · neu in diesem Jahr · P10 · baut auf: $m$ und $n$ ablesen (Einheit 2), Punktprobe und Nullstelle (Einheit 3), lineare Gleichungen, Terme zusammenfassen}
\erg{32}{a) $m = 1$, $n = -3$, $y = x - 3$ \quad b) $m = -3$, $n = 2$, $y = -3x + 2$ \quad c) $m = 2$, $n = 4$, $y = 2x + 4$}
\erg{33}{a) $n = 3$, $y = 2x + 3$ \quad b) $n = -2$, $y = 3x - 2$ \quad c) $n = 4$, $y = -x + 4$ \quad d) $4 = 6 + n$, $n = -2$, $y = -2x - 2$ \quad e) $3 = -2 + n$, $n = 5$, $y = 0{,}5x + 5$}
\erg{34}{a) $m = 3$, $n = 2$, $y = 3x + 2$ \quad b) $m = 2$, $n = -1$, $y = 2x - 1$ \quad c) $m = 9 : 3 = 3$, $n = 3$, $y = 3x + 3$ \quad d) $m = -9 : 3 = -3$, $n = 4$, $y = -3x + 4$}
\erg{35}{a) Gerade durch $(0|1)$ und $(2|5)$ \quad b) Gerade durch $(-3|4)$ und $(3|{-2})$ \quad c) $m = (1 - 4) : (4 - (-2)) = -3 : 6 = -0{,}5$; $4 = -0{,}5 \cdot (-2) + n$, $n = 3$; $y = -0{,}5x + 3$}
\erg{36}{a) $2x = 4$, $x = 2$, $S(2|8)$ \quad b) $2x = 8$, $x = 4$, $S(4|13)$ \quad c) $3x = 6$, $x = 2$, $S(2|6)$ \quad d) $-3x = -9$, $x = 3$, $S(3|13)$}
\erg{37}{e) $-3x = -9$, $x = 3$, $S(3|4)$ \quad f) $1{,}5x = 3$, $x = 2$, $S(2|2)$ \quad g) $-4x = -6$, $x = 1{,}5$, $S(1{,}5|0)$}
\erg{38}{a) Mia hat oben $5 - 11$ gerechnet, unten aber $4 - 2$: die Reihenfolge ist vertauscht. Richtig: $m = (11 - 5) : (4 - 2) = 3$; $5 = 3 \cdot 2 + n$, $n = -1$; $y = 3x - 1$ \quad b) $m = (7 - (-1)) : (3 - 1) = 4$, $n = -5$, $y = 4x - 5$}
\erg{39}{a) parallel: gleiches $m = 2$, verschiedenes $n$ \quad b) schneiden sich: verschiedene Steigungen ($-1$ und $1$), und zwar in $(0|2)$ \quad c) liegen aufeinander: gleiches $m = 0{,}5$ und gleiches $n = 1$}
\erg{40}{a) $m = (700 - 820) : (10 - 4) = -20$; $820 = -20 \cdot 4 + n$, $n = 900$; $y = -20x + 900$ \quad b) Der Ballon sinkt 20 m je Minute; zu Beginn der Messung ist er 900 m hoch. \quad c) $0 = -20x + 900$, $x = 45$: nach 45 Minuten}
